\subsection{実際の最終形態(卒論=pdf)への変換}
hiki+keynoteで卒論の内容ができたら，それを卒論の最終形態つまりpdfへ変換する必要が出てきます．

ここで問題が発生します．hikiよりもlatexの方が機能が豊富なこと．
つまり，簡易に書くにはhikiなどのmark downで書いていくのがいいんですが，
文書として完成度を高めるには，latexで細かい設定を調整する必要がどうしても出てきます．

具体的には，
\begin{enumerate}
\item 図の配置を調整するwrapの数値調整
\item 参照文献の記述
\item リスト，引用の体裁
\item 章立て階層構造
\end{enumerate}
などが問題になるところです．

これはhikiではどうしようもありません．一つの手はもう少し高機能なmark up言語
例えばasciidocなどに変更することですが，これはどこまでいっても終わりのない
方向のようで，結局はlatexで書いているのと同じになる可能性があります．

我々は違う戦略をとります．それは，
\begin{quote}\begin{verbatim}
latexをベースにして，hikiを生成する
\end{verbatim}\end{quote}
と言う手です．

DRY(Don't Repeat Yourself)原則さえ守れば，文書管理はいいのですから，
ある段階まですすめばhikiではなくlatexを原本にするのです．
そのための変換器latex2hikiとその派生rake環境が用意されています．
そちらはこちらで卒業後に変換します．参照文献はlatexに移してから
修正してください．

ここでは，hikiでできることをとことん突き詰めておきます．

\subsubsection{install}\begin{quote}\begin{verbatim}
hiki -i
\end{verbatim}\end{quote}
でinstallされています．新たに使うコマンド群は次の通りです．
\begin{lstlisting}[style=]
rake change_wrap     # change latex figures to wrap format
rake latex           # latex conversion FILE1
rake latex_all       # latex conversion whole hiki files in the current dir
rake latex_base      # latex conversion FILE1(hiki) to FILE2(latex)
rake latex_wrap      # latex conversion FILE1 with wrap format
\end{lstlisting}\begin{quote}\begin{verbatim}
rake latex hikiutils_bob.hiki
\end{verbatim}\end{quote}
と打ち込むと

\subsubsection{rake latex手順（個別ファイルの変換）}\begin{quote}\begin{verbatim}
rake latex sync.hiki
\end{verbatim}\end{quote}
とすると，
\begin{quote}\begin{verbatim}
latex_dir/sync.tex
\end{verbatim}\end{quote}
にlatex変換後のファイルが生成します．これを，TeXShopでcommand\_t変換します．

うまくいかないときは，terminalで
\begin{quote}\begin{verbatim}
platex sync.tex
dvipdfmx sync.dvi
\end{verbatim}\end{quote}
を試してみてください．

\paragraph{注意}
hikiの初めの部分は，
\begin{quote}\begin{verbatim}
!title:hikiutils -iによる卒論作成システム
!author:Shigeto R. Nishitani
!date: Kwansei Gakuen Univ., 2017/1
\end{verbatim}\end{quote}
とすると，
\begin{lstlisting}[style=]
\begin{document}
\title{hikiutils -iによる卒論作成システム}
\author{Shigeto R. Nishitani}
\date{ Kwansei Gakuen Univ., 2017/1}
\maketitle
\end{lstlisting}
と変換してくれます．

必要な図は，figsから自動で取るようになっていますが，サイズは解像度が問題のときは手動で調整してください．

\subsubsection{rake latex\_all手順（ディレクトリーないの一括変換）}\begin{quote}\begin{verbatim}
bob% cat hikiutils_bob.hiki 
!hikiutilsを用いた卒業論文作成

! [[hikiutils_bob_sync]]
! [[hikiutils_bob_latex_all]]

\end{verbatim}\end{quote}
とpwdのdirectoryの同名のbasenameに.hikiの拡張子がついたファイルが用意されていて，
\begin{quote}\begin{verbatim}
rake latex_all
\end{verbatim}\end{quote}
をおこなうと自動で全部を一体にまとめた文書へのlatex変換が出来上がります．
開いているhikiutils\_bob.texをTeXShopで変換してみてください．

うまくできないときは\verb|ここ(http://qiita.com/hideaki_polisci/items/3afd204449c6cdd995c9#51-texshopの設定)|を参照して自力で入れてみてください．だめならdonkey.

\begin{figure}[htbp]\begin{center}
\includegraphics[width=6cm,bb=0 0 442 500]{../figs/./hikiutils_bob.pdf}
\caption{}
\label{default}\end{center}\end{figure}
というようになります．

\paragraph{下準備}
latex\_dir内に幾つかのtex雛形を入れておく必要があります．
\begin{quote}\begin{verbatim}
hiki -i
\end{verbatim}\end{quote}
とhikiutils環境を再度初期化すると，自動でインストールする設定です．
なかったら手動で作ってください．

\paragraph{head.tex}\begin{lstlisting}[style=]
bob% cat head.tex
\title{卒業論文\\
\vspace{4cm} hikiutilsを用いた\\卒業論文作成}
\author{ 関西学院大学　理工学部　情報科学科\\\\1234 西谷滋人}
\date{\vspace{3cm} 2017年  3月\\
\vspace{3cm} 指導教員　　西谷　滋人　教授}
\maketitle

\end{lstlisting}
\paragraph{pre.tex}\begin{lstlisting}[style=]
bob% cat pre.tex
\documentclass[10pt,a4j]{article}

\def\Vec#1{\mbox{\boldmath $#1$}}
\usepackage[dvipdfmx]{graphicx}

\setlength{\textheight}{275mm}
\headheight 5mm
\topmargin -20mm
\textwidth 160mm
\textheight 250mm
\oddsidemargin -0mm
\evensidemargin -5mm

\pagestyle{empty}
\makeatletter
  \def\@maketitle{%
  \newpage\null
  \vskip 2em%
  \begin{center}%
  \let\footnote\thanks
    {\large\bf \@title \par}%
    \vskip 1.5em%
    {\large\bf \@author \par}%
    \vskip 1.5em%
    {\small \@date}%
  \end{center}%
}
\makeatother

%\documentclass[10pt, a4j]{article}
%%\usepackage{citesort}
\usepackage{amssymb}
\usepackage[dvipdfmx]{graphicx}% 図を入れるときに使用
\usepackage{wrapfig}% 図の周りに本文を流し込みたいときに使用
\usepackage{subfigure}
\usepackage{here}
\end{lstlisting}
